% Root, B and D: full actual-file and exact primary-source audits passed for all three statements.
\begin{theorem}[Proved: constant second cohomology of a complete local surface puncture]
\label{D:local-surface-constant-H2}
Let $(R,\mathfrak m,k)$ be a complete excellent normal local
domain of dimension two, with algebraically closed residue field.
Put $U=\operatorname{Spec}R-\{\mathfrak m\}$ and
$G=\pi_1(U)$. There is an integer $b_R$ such that, for every
positive integer $d$ invertible in $R$,
\[
 |H^2_{\mathrm{cont}}(G,\mathbf Z/d\mathbf Z)|
 \le |H^2_{\mathrm{et}}(U,\mathbf Z/d\mathbf Z)|
 \le d^{b_R}.                                         \tag{D.LH2.1}
\]
The group action on the displayed coefficients is trivial.
In particular these groups are finite. The same bound, with
$d^{b_R}$ replaced by $|A|^{b_R}$, holds for every finite
constant abelian group $A$ whose order is invertible in $R$.
No assertion about arbitrary nonconstant coefficient systems
is needed here.
\end{theorem}
\begin{proof}
The puncture $U$ is regular. By resolution of excellent normal
surfaces followed by embedded resolution, choose a proper
birational morphism
\[
             \pi:Y\longrightarrow\operatorname{Spec}R
\]
which is an isomorphism over $U$, with $Y$ regular and
$E=(\pi^{-1}\{\mathfrak m\})_{\mathrm{red}}$ a strict normal
crossings divisor. If necessary, a blowup of the closed point
first makes the exceptional fibre nonempty. The precise
resolution inputs are Theorem~54.14.5 and Lemma~54.15.6
\citep[Tags 0BGP and 0BIC]{StacksDSurfaceResolution}.
The components $E_1,\ldots,E_v$ are smooth proper curves over
$k$. Let $g_i$ be their genera, and let $e$ be the number of
their intersection points, counting distinct intersection
points separately. There are no triple intersections.

Set $\Lambda=\mathbf Z/d\mathbf Z$ and let $j:U\hookrightarrow Y$.
Gabber's absolute purity theorem, in its regular normal-crossings
form, gives
\[
 \begin{split}
 R^0j_*\Lambda&=\Lambda_Y,\\
 R^1j_*\Lambda&=\bigoplus_i\Lambda_{E_i}(-1),\\
 R^2j_*\Lambda&=\bigoplus_{x\in\operatorname{Nodes}(E)}
                                      \Lambda_x(-2),\\
 R^qj_*\Lambda&=0\quad(q>2).
 \end{split}                                         \tag{D.LH2.2}
\]
Here a sheaf on a closed stratum denotes its direct image to
$Y$. We use exactly Expos\'e~XVI, Theorem~3.1.1 and
Corollary~3.1.3 of \citet{A_ILO2014}; the complete statements
and the proof of the normal-crossings corollary appear on
pp.~242--243 of the preserved arXiv version. That corollary
applies to regular $\mathbf Z[1/d]$-schemes. It does not
require finite type over a field. In particular its hypotheses
apply to this excellent resolution. This is a use of purity
on the regular pair $(Y,E)$, not a constructibility assertion
for an arbitrary open immersion into the singular local scheme.

The base is strictly henselian. Global sections on its small
etale topos are exact, and proper base change consequently gives
\[
                   H^a(Y,\Lambda)=H^a(E,\Lambda).
\]
For the proper base-change input we use
\citep[Theorems~17.7 and 17.10]{MilneLEC2013H2D}.
Nilpotents in the closed fibre do not change its etale topos.
Normalization of the nodal curve $E$ and the node exact
sequence give $|H^2(E,\Lambda)|=d^v$; for each smooth
component $|H^1(E_i,\Lambda)|=d^{2g_i}$.
Tate twists are constant here, since the algebraically closed
residue field contains all $d$th roots of unity and Hensel's
lemma lifts them to $R$.

Apply the Leray spectral sequence
\[
 H^a(Y,R^qj_*\Lambda)\Longrightarrow H^{a+q}(U,\Lambda).
\]
The three terms of total degree two in its second page have
orders $d^v$, $d^{2\sum_i g_i}$ and $d^e$, respectively,
by (D.LH2.2). The corresponding limiting terms are subquotients
of these finite groups. Thus
\[
                |H^2(U,\Lambda)|
                    \le d^{v+2\sum_i g_i+e}.
\]
We may take $b_R=v+2\sum_i g_i+e$. Equivalently, one can obtain
finiteness from the localization triangle and the finite
Mayer--Vietoris complex of supports in the $E_i$, using purity
for the divisors and their intersection points.

Finally the usual degree-two comparison is injective:
\[
                H^2(G,\Lambda)\lhook\joinrel\longrightarrow
                                H^2_{\mathrm{et}}(U,\Lambda).
\]
Indeed the Cartan--Leray spectral sequence for all connected
finite Galois covers has degree-zero geometric term $\Lambda$,
and degree-one geometric term zero: every finite-coefficient
degree-one class is a finite torsor and is killed by a further
finite etale cover. The only possible incoming differential
to $(2,0)$ is therefore zero, and this limiting term is the
last filtration piece in total degree two. This is the same
comparison used in
Theorem~\ref{D:surface-tame-cross-characteristic-H2}; see
\citep[Theorem~14.9, Remark~14.10 and Proposition~14.13]
{MilneLEC2013H2D}. Decomposing a finite constant abelian
group into cyclic groups proves the final assertion.
\end{proof}

\begin{corollary}[Proved: one subgroup kills all bounded local central classes]
\label{D:local-bounded-central-obstructions}
In the preceding setting, for every fixed integer $B$ there
is an open normal subgroup $G_B\subset G$ such that
\[
 H^2(G,\mathbf Z/d\mathbf Z)\longrightarrow
 H^2(G_B,\mathbf Z/d\mathbf Z)
\]
is zero for every $d\le B$ invertible in $R$.
The subgroup is independent of the representation rank and
coefficient extension degree; no bound uniform in $B$ is asserted.
\end{corollary}
\begin{proof}
There are finitely many relevant $d$ and, by the theorem,
finitely many classes for each. Every normalized continuous
two-cocycle with finite coefficients factors through one
finite quotient of $G$: compactness refines its finitely
many clopen fibres by rectangles of cosets of a common open
normal subgroup. Its restriction to that subgroup is the
zero cocycle. Intersect these finitely many kernels.
Equivalently, use the continuous-cochain comparison
\citep[II, Proposition~4.2]{MilneCFT2020D}.
\end{proof}


\begin{corollary}[Proved: prime-to-characteristic cohomology on arbitrary-dimensional strict-local opens]
\label{D:strict-local-open-cohomology-finite}
Let $R$ be a strictly henselian quasi-excellent Noetherian local
ring, let $j:U\hookrightarrow\operatorname{Spec}R$ be an open
immersion, and let $d$ be invertible in $R$. For every constructible
sheaf $\mathcal F$ of $\mathbf Z/d\mathbf Z$-modules on $U$,
all groups $H^q(U,\mathcal F)$ are finite. In particular, if $U$
is nonempty and connected, then $H^2(\pi_1U,\mathbf Z/d\mathbf Z)$ is finite.
For each fixed $B$ one open normal subgroup of $\pi_1U$ kills
all constant central classes with $d\le B$ invertible in $R$.
These assertions impose no dimension-two hypothesis. They do
not assert a uniform exponent in $d$ analogous to (D.LH2.1).
\end{corollary}
\begin{proof}
Since $R$ is Noetherian, its open subscheme $U$ is quasi-compact,
so $j$ is of finite type. Expos\'e~XIII, Theorem~1.1.1(i) of
\citet{A_ILO2014} says exactly that, for a finite-type morphism
into a quasi-excellent Noetherian scheme and constructible
$\mathbf Z/d\mathbf Z$ coefficients with $d$ invertible on the
base, every higher direct image is constructible. Apply this
to $j$. In particular each $R^qj_*\mathcal F$ has a finite
stalk at the closed geometric point.

For a strictly local scheme, taking global sections on the
small etale topos is the closed geometric-stalk functor, hence
is exact. Indeed every etale neighbourhood of that point
has a section through it, so both equality of sections and
existence of local lifts can be tested after pulling back
along such a section. The Leray spectral sequence consequently
identifies
\[
 H^q(U,\mathcal F)
   =\Gamma(\operatorname{Spec}R,R^qj_*\mathcal F)
   =(R^qj_*\mathcal F)_{\overline{\mathfrak m}},
\]
which is finite. The same Cartan--Leray degree-two injection
as above gives the fundamental-group assertion. Finally apply
the finite-cocycle argument of
Corollary~\ref{D:local-bounded-central-obstructions} to the
finitely many groups and classes of bounded order.
\end{proof}
